\section{Simulations}
\label{sec:simulation}

We check the theory with the estimator a user would run, not only with
population calculations. Both experiments generate multi-site data
from the fixed-effects model \eqref{eq:response-model}, estimate the
five bound inputs with the rank consistent estimator of
\citet{boyarsky2026}, and solve the estimated bound programs. The
replication code reproduces every number and figure.

\subsection{Design}
\label{subsec:sim-design}

The context space has $p=6$ coordinates. Six existing sites carry fixed
profiles whose centered span $\Vspan$ has rank three, so three
directions of the fixed effects are unidentified, the same mechanism as
in the microcredit data. Each site contributes 500 observations with
three unit-level covariates, a binary covariate, a discrete covariate
whose distribution varies by site, and a Gaussian covariate. Treatment
probabilities are known and vary by site between $0.35$ and $0.65$. An
independent target covariate sample has 1{,}500 observations. The
target region $\targset$ holds fifteen fixed profiles, five inside the
affine hull of the pool and ten with components outside it. The
restriction is the per-arm box $|\beta_{a,j}|\leq 2.5$, which keeps
every width finite and contains the truth with room to spare. Estimation
uses four folds and correctly specified linear and multinomial logistic
nuisance fits, so the experiments isolate the selection problem rather
than nuisance error. Exact population moments and widths come from
finite enumeration of the discrete covariates. Each experiment runs 40
replicates that redraw the data with the geometry fixed.

\subsection{Selecting among four candidates}
\label{subsec:sim-selection}

Four candidate sites cover the three categories of \cref{sec:selection}: $k_1$
inside the convex hull of the existing profiles, $k_2$ inside the affine
hull but outside the convex hull, and $k_3$ and $k_4$ outside the
affine hull along two different directions, $k_3$ far outside along the
first and $k_4$ a short step along the second. For each candidate we
compute the design-time prediction from the existing pool alone, then
simulate the candidate experiment from the same model, refit the
estimator on the seven-site pool, and record the realized worst-case
width over $\targset$.

\Cref{fig:selection} reports the results. The population worst-case
width is $18.31$ for the existing pool and stays exactly $18.31$ when
$k_1$ or $k_2$ is added: sites inside the affine hull change nothing
about identification, whether or not they are inside the convex hull.
The two hull-expanding candidates cut the population worst-case width to
$14.63$ ($k_3$) and $13.54$ ($k_4$), reductions of 20 and 26 percent.
The estimated widths track the population values to within sampling
noise in every cell, with mean estimates of $18.31$, $18.31$, $18.29$,
$14.63$, and $13.59$. The estimated selection picks $k_4$, the
population-optimal candidate, in all 40 replicates, and so does the
design-time selection computed before the candidate data exist. The
design-time predictions are conservative, $16.49$ against a realized
$14.63$ for $k_3$ and $14.04$ against $13.59$ for $k_4$, because the
neutral forecast slices the restriction box through its widest
cross-section, and they preserve the ranking.

The comparison between $k_3$ and $k_4$ makes the geometry of
\cref{sec:geometry} concrete. The nearby candidate beats the distant
one because the worst-case targets load more on the direction $k_4$
resolves than on the direction $k_3$ resolves. Distance to the hull does not enter the identified set. The identified
set depends on which direction a candidate resolves and how much the
target region loads on it.

\begin{figure}[t]
  \centering
  \includegraphics[width=\textwidth]{figure1.png}
  \caption{\textbf{Site selection with estimated bounds.} Left: the
  fixed geometry, projected onto one direction inside the existing span
  and the first direction outside it. Right: worst-case width over the
  target region for the existing pool and after each candidate, showing
  the mean and standard deviation of the estimated widths across 40
  replicates, the exact population width, and the design-time
  prediction computed before the candidate experiment.}
  \label{fig:selection}
\end{figure}

\subsection{The conflict mechanism}
\label{subsec:sim-monotonicity}

The second experiment tests \cref{thm:monotone} and its failure mode.
Two profiles are reused, the in-span candidate $k_2$ and the
hull-expanding candidate $k_3$, and each is added twice: once with
outcomes drawn from the maintained model, labels $a_1$ and $a_3$, and
once with the candidate site's treated-arm mean shifted by a conflict
of size $0.75$ whose sign is drawn at random each replicate, labels
$a_2$ and $a_4$. The restriction box is re-centered at random each
replicate so that it is generically asymmetric around the identified
slice, which is when the translation mechanism of
\cref{sec:monotonicity} can move widths. Population widths and truth
coverage are computed exactly for every drawn box.

\Cref{fig:monotonicity} and its underlying table say three things.
First, the model-consistent additions behave exactly as
\cref{thm:monotone} requires: $a_1$ leaves the population worst-case
width unchanged in all 40 replicates, $a_3$ strictly shrinks it in all
40, with reductions between $1.1$ and $6.7$, and neither ever loses the
truth. Second, the conflicted in-span addition $a_2$ moves the
population worst-case width in both directions, widening it in 45
percent of replicates and narrowing it in 55 percent, by up to $0.71$
either way. The direction is unpredictable in advance, which is the
point of the counterexample. Third, and most important, $a_2$ removes the true target effect from the
population identified set at the span-identified targets in every
replicate, while its worst-case width moves by at most a few percent.
The width does not show the damage, and the interval no longer covers
the true effect. The conflicted hull-expanding addition
$a_4$ still shrinks widths here, and its bias lies in the directions that remain unidentified. This is
why the compatibility screen of \cref{sec:monotonicity}, and not the
width, is the right check for it.

\begin{figure}[t]
  \centering
  \includegraphics[width=\textwidth]{figure2.png}
  \caption{\textbf{What a conflicting site does.} (A) The two-site
  example of \cref{sec:monotonicity}: a third site on the pool's own
  line whose outcomes conflict with the pool moves the identified slice
  and doubles the width at the target. (B) Mean estimated worst-case
  width across 40 replicates with the exact population value marked.
  (C) Post-addition against existing-pool worst-case width by
  replicate. Open markers are replicates whose population identified
  set no longer covers the true effect at every target: all 40
  replicates of the conflicted in-span addition $a_2$, and none of the
  model-consistent additions.}
  \label{fig:monotonicity}
\end{figure}

\FloatBarrier
